\section{Выводы}

В лабораторной работе был реализован алгоритм Ахо--Корасик для поиска большого числа образцов в тексте
над числовым алфавитом. Основная сложность заключалась не в построении trie, а в корректном добавлении к
нему связей неудачи и выходных ссылок: именно они позволяют не начинать поиск заново после несовпадения
и находить несколько образцов, заканчивающихся в одной позиции.

Важной частью реализации стало соблюдение варианта алфавита. Переходы автомата строятся по значениям
\texttt{uint32\_t}, а не по символам строкового представления числа. Поэтому записи с ведущими нулями
обрабатываются правильно, а реализация не сводит задачу к поиску по меньшему алфавиту.

Программа обрабатывает текст потоково и не хранит его целиком в памяти. Для восстановления позиции начала
вхождения достаточно очереди последних позиций длиной не больше максимальной длины образца. Это делает
решение применимым к входам без ограничений на длину строк и количество чисел.

Итоговая реализация соответствует условию задачи: строит автомат по множеству образцов, выполняет один
проход по тексту и выводит для каждого вхождения номер строки, номер слова и номер образца.
\pagebreak
